% Part of main.tex; compile main.tex.
% ---------------------------------------------------------------------------
\section{Introduction}
\label{sec:intro}

Many high-performance computing (HPC) simulations and analysis codes
alternate between computation and Input/Output (I/O) phases. A simulation
advances a number of timesteps and then writes a checkpoint or plot file,
while an analysis code reads a timestep and then processes it. When I/O is
performed synchronously, the computation stalls until the I/O phase
completes. Because storage bandwidth has grown more slowly than compute
capability, these stalls can become a significant fraction of the total run
time.

Asynchronous I/O addresses this problem by separating the issue of an I/O
operation from its completion. The application initiates an I/O operation,
continues with its next computation phase, and only waits for completion
when it needs the buffer again or reaches a synchronization point.
Implementations range from POSIX asynchronous I/O (AIO) and non-blocking
MPI-IO to background-thread frameworks, such as the HDF5 asynchronous
Virtual Object Layer (VOL) connector~\cite{tang2019async,tang2022async}, and
server-based approaches that use dedicated cores or staging
resources~\cite{dorier2016damaris}. In the ideal case, the observed I/O cost
per step falls from $t_{io}$ to a small overhead, and the time per step
becomes approximately $\max(t_{comp}, t_{io})$ rather than
$t_{comp} + t_{io}$~\cite{ravi2023evaluating}.

Whether this ideal is reached depends on a question that is easy to
overlook: \emph{what is the application doing while the I/O runs in the
background?} The background I/O engine needs CPU time, memory bandwidth for
buffer copies, and network bandwidth to reach the parallel file system (PFS).
A real computation phase competes for all three. However, a common practice
in the asynchronous I/O literature is to replace the computation phase with
a call to \texttt{sleep()} of a chosen
duration~\cite{li2021h5bench,tang2022async,ravi2023evaluating}. This choice
is convenient, because it makes experiments cheap and allows the authors to
select a compute phase long enough to fully hide the I/O. However, it also
removes the contention that could prevent the overlap.

Contention can appear on either side of the overlap. The background engine
can take CPU time, memory bandwidth, or MPI resources from the application,
which makes the computation run longer. Alternatively, the application can
starve the engine, so the I/O does not progress during the computation and
the application waits for it at the next synchronization point. Studies that
report only the observed I/O time cannot clearly see either effect, and
studies that use \texttt{sleep} cannot see them at all, because a sleeping
thread neither competes for resources nor withholds them. For this reason,
we evaluate asynchronous I/O by the quantity the application actually
experiences, its end-to-end run time, and attribute any difference to the
compute side or the I/O side.

In this paper, our contributions are the following:
\begin{itemize}
  \item We survey how 17 asynchronous I/O and overlap evaluations spanning
        HDF5, ADIOS, PDC, POSIX-level checkpointing, and deep-learning I/O
        treat the computation phase. Only 6 of them use a computation that
        consumes the resources background I/O needs
        (Section~\ref{sec:survey}).
  \item We collect the evidence that bears on the validity of
        \texttt{sleep}-based evaluation from interference studies,
        real-application evaluations, overlap benchmarks, and the caveats
        stated by the \texttt{sleep}-based studies themselves.
  \item We present an open benchmark that drives POSIX, MPI-IO, HDF5, and
        ADIOS2 with interchangeable compute surrogates calibrated to the same
        standalone duration. We use it to measure how far
        \texttt{sleep}-based estimates of the benefit of asynchronous I/O
        diverge from estimates that include contention
        (Section~\ref{sec:eval}).
  \item We provide recommendations for designing future evaluations and for
        interpreting published results that rely on \texttt{sleep()}.
\end{itemize}

The remainder of the paper is organized as follows.
Section~\ref{sec:related} reviews related work. Section~\ref{sec:survey}
describes our survey methodology and its findings. Section~\ref{sec:eval}
presents our benchmark and experimental results. Lastly,
Section~\ref{sec:conclusion} presents our recommendations and concludes the
paper.
